% ECONOMICS_PROBLEM: {"id": "E62-543", "title": "State-dependent fiscal multipliers", "jel_code": "E62", "source_id": "OP-0251"}
% !TeX program = xelatex
\documentclass[11pt,a4paper]{article}
\usepackage[margin=23mm]{geometry}
\usepackage{fontspec}
\setmainfont{Latin Modern Roman}
\usepackage{amsmath,amssymb,mathtools}
\usepackage{xcolor,enumitem,booktabs,longtable,array,xurl,hyperref}
\definecolor{muted}{HTML}{53636C}
\hypersetup{colorlinks=true,urlcolor=blue,linkcolor=blue}
\setlength{\parindent}{0pt}
\setlength{\parskip}{5pt}
\newcommand{\R}{\mathbb R}
\newcommand{\E}{\mathbb E}
\newcommand{\N}{\mathbb N}
\newcommand{\ind}{\mathbf 1}
\newcommand{\Var}{\operatorname{Var}}
\newcommand{\Cov}{\operatorname{Cov}}
\newcommand{\supp}{\operatorname{supp}}
\DeclareMathOperator*{\argmin}{arg\,min}
\DeclareMathOperator*{\argmax}{arg\,max}
\newcommand{\entrymeta}[3]{{\sffamily\small\color{muted}\textbf{\texttt{#1}}\quad|\quad #2\par #3\par}}
\newcommand{\Problemref}[1]{\texttt{#1}}
\newcommand{\JELref}[2]{\texttt{#2} (JEL \texttt{#1})}
\begin{document}
\section*{State-dependent fiscal multipliers}
\textbf{Problem ID:} \texttt{E62-543}\quad \textbf{JEL:} \texttt{E62}\par
% BEGIN ECONOMICS STATEMENT
\label{problem:OP-0251}
{\sffamily\small\color{muted}Original subject group: Fiscal policy and sovereign debt\par}
\label{definitions:problem:30007532}
\textbf{Audit classification:} Explicit published open question. Size/sign multiplier asymmetries are expressly debated; this causal surface also adds distinct state variation. Verification concerns the cited version, not first priority or proof that this exact editorial specification remains unsolved on October 8, 2026.
\subsection*{Definitions and precise research target}
\paragraph{Editorial mathematical specification.} Let \(G_t\) be real government purchases and \(Y_t\) real output. Define predetermined state \(s_t\) using a predeclared recession indicator, inflation, debt, openness, and monetary-policy constraints. An instrument \(Z_t\) identifies an unanticipated purchase innovation, excluding responses to anticipated future output and other simultaneous fiscal measures. For intervention magnitude \(a\), let \(Y_{t+h}(a)\) and \(G_{t+h}(a)\) be potential outcomes under a stated financing and monetary-policy rule. Define the cumulative state-dependent multiplier
\[
m_H(s,a)=
\frac{\sum_{h=0}^H E[Y_{t+h}(a)-Y_{t+h}(0)\mid s_t=s]}
{\sum_{h=0}^H E[G_{t+h}(a)-G_{t+h}(0)\mid s_t=s]},
\]
when its denominator is nonzero. Identify the response surface for a fixed horizon and observed range of shock sizes, with confidence sets robust to weak instruments and overlapping outcomes. Determine which state differences remain after holding instrument composition, financing, and policy reactions comparable. Distinguish marginal from large-intervention multipliers and positive from negative spending shocks; linear local projections cannot automatically establish those equalities. Report support limitations when large shocks occur only in exceptional crises. The goal is a transportable conditional estimate under explicit identification assumptions, not one invariant fiscal multiplier applicable to every economy and spending program.

Purchases and output must be real flows in the same currency and period units; scaling a policy-rate shock or a tax change is not the same intervention. Specify the full intervention path \(\delta G_{t:t+H}(a)\), fiscal financing, and monetary reaction, so that equal \(a\) means comparable spending across states. In a derivative multiplier replace numerator and denominator by their derivatives at \(a=0\), with a nonzero denominator. A just-identified instrumental-variable ratio generally identifies an instrument-specific local response and does not identify the full nonlinear surface \(m_H(s,a)\) without richer support and exclusion assumptions. Predetermined recession status must be based on information available before the shock, not output subsequently changed by it. With continuously indexed states, identification is defined only on their population support.
\subsection*{Variants and assumption changes}
The following are \emph{editorial variants}, not additional published open conjectures.
\begin{enumerate}
\item \emph{Monetary constraint.} Compare the same purchase and financing path under a predetermined effective-lower-bound regime and an unconstrained policy regime. Specify selection into the regime and sufficient instrument support before interpreting a multiplier difference causally.
\item \emph{Spending composition.} Split purchases into public investment and government consumption with separately identified innovations. Hold total spending and financing constant and estimate composition effects, rather than pooling instruments that change different programs.
\end{enumerate}
\subsection*{References and source audit}
\begin{itemize}
\item Adrien Auclert; Matthew Rognlie; Ludwig Straub. \emph{Fiscal and Monetary Policy with Heterogeneous Agents}. 2025. Locator: Section5, Nonlinearities, printed pp.19--20; size/sign asymmetries rather than all state-dependence questions. Primary text checked. \url{https://web.stanford.edu/~aauclert/annual_review_hank.pdf}.
\end{itemize}
Size/sign multiplier asymmetries are expressly debated; this causal surface also adds distinct state variation.
\begin{enumerate}\raggedright
\item\hypertarget{definitions:source:30007532:1}{} Adrien Auclert; Matthew Rognlie; Ludwig Straub. 2025. \emph{Fiscal and Monetary Policy with Heterogeneous Agents}. Section5, Nonlinearities, printed pp.19--20; size/sign asymmetries rather than all state-dependence questions.
\url{https://web.stanford.edu/~aauclert/annual_review_hank.pdf}.
DOI: \url{https://doi.org/10.1146/annurev-economics-091624-044646}.
\end{enumerate}

\subsection*{Common conventions: Source mathematical conventions}

These conventions apply to each entry unless explicitly replaced.

\textbf{Models and identification.} A primitive model \(m\) specifies all probability laws, feasible choices, information, equilibrium and measurement rules used by its target. The admissible class \(\mathcal M\) is fixed independently of outcomes. If \(P_O(m)\) is its observable probability law and \(\tau(m)\) the target, its identified set at observable law \(P\) is
\[
 \mathcal I_\tau(P)=\{\tau(m):m\in\mathcal M,\ P_O(m)=P\}.
\]
Point identification means that this set is a singleton. An empty set signals incompatibility of data and assumptions. Coordinate bounds need not describe the joint set. Bounds are sharp only if every retained target is attainable or arbitrarily approachable by compatible models. Finite bounds require explicit restrictions on unobserved quantities; unrestricted confounding, hidden wealth, extrapolation or arbitrary equilibrium selection can yield uninformative sets.

\textbf{Causal interventions.} A potential outcome \(Y_i(p)\) is the outcome under a complete policy regime \(p\), including timing, financing, information and market spillovers. The target population and units are fixed before treatment. With interference, use the complete assignment vector or a justified exposure mapping; individual no-interference notation alone cannot identify an economy-wide intervention. A difference of mean potential outcomes does not require the cross-world joint distribution. Conditioning on post-treatment migration, survival, enrollment or employment changes the estimand and needs separate justification.

\textbf{Statistical statements.} An estimator, confidence set, forecast or test specifies a sampling law, model class and loss/coverage criterion. Uniform \((1-\alpha)\)-coverage means \(\inf_{m\in\mathcal M}\Pr_m\{\tau(m)\in C_n\}\geq1-\alpha\), or an explicitly asymptotic version. The whole parameter space trivially has coverage, so informativeness requires a stated width, risk or power objective. A forecasting improvement is relative to a named benchmark, information set and evaluation population.

\textbf{Equilibrium and optimization.} An equilibrium comprises feasible plans, optimality, beliefs where needed, budget constraints and all market/resource clearing. The primitives or a stated benchmark define each correspondence \(\mathcal E(p;m)\). Nonemptiness is assumed or proved, never inferred just from compact policy sets. A selected-equilibrium target uses a declared selection rule; a robust target quantifies over all permitted equilibria. Use suprema or infima unless attainment is established. An \(\varepsilon\)-optimal policy is feasible and within \(\varepsilon\) of the finite optimum value. Report an empty feasible set separately.

\textbf{Welfare and accounting.} Normative utility, interpersonal normalization, social weights, numeraire, discounting and the use of public receipts are primitives. Behavioral utility need not coincide with the welfare criterion. Household welfare includes ownership income and rents once; adding profits again to compensating variation generally double-counts them. A monetary equivalent is defined by an explicit transfer and price convention, with monotonicity and range conditions. Average effects, mechanism contributions, transfers and welfare losses are different targets. Infinite sums require convergence and dynamic feasible plans require terminal or no-Ponzi conditions.

\textbf{Source versions.} A working-paper date, revision date and journal publication year are distinguished. A DOI for a journal version is not automatically the DOI of an earlier manuscript. Locators refer to the stated version and printed pages where specified. A metadata-only check does not verify a claim about a full-text conclusion. Every entry's source subsection narrows the provenance claim accordingly.
% END ECONOMICS STATEMENT
\end{document}
